\subsection{A hardest completion and randomized adaptive optimality}
\label{sec:adaptive-dual-completion}

Theorem~\ref{prop:risk_classification} already gives equality of randomized
batch and adaptive mean-width risks for this model. The specialization below
extracts the common hardest completion directly from the finite-game dual,
providing a computable optimality certificate for the compact solver.
Adaptive policies retain the same hard budget $B$, independent random seed,
exact coordinate information and sharp terminal interval as in
Section~\ref{sec:risk_setup}.

\begin{proposition}[A common hardest completion]
\label{prop:adaptive-dual-completion}
For the disjoint exact-total model, an optimum $w^*$ of
\eqref{eq:random-dual} defines the active-block values
$y^*_{gi}=w^*_{gi}$ of a feasible residual completion. Restore each omitted
constant-coefficient block with any fixed nonnegative allocation of its
total, and restore the determined singleton and zero-total values.
For every admissible terminal coordinate set $Q$,
\begin{equation}
 W(Q,y^*)\geq R_B.
 \label{eq:common-hardest-width}
\end{equation}
Consequently, if $\mathcal A_B$ is the class of randomized adaptive
coordinate-query policies satisfying the hard budget, then
\begin{equation}
 \inf_{A\in\mathcal A_B}\sup_{y\in\F}
       \mathbb E[W_A(y)] = R_B.
 \label{eq:adaptive-random-equality}
\end{equation}
The existing compact randomized batch policy attains this value. At $y^*$,
every execution of every admissible adaptive policy has terminal width at
least $R_B$, not merely an expectation bounded below by $R_B$.
\end{proposition}

\begin{proof}
The dual equalities $\sum_i w^*_{gi}=E_g$ and inequalities $w^*_{gi}\geq0$
are exactly the active residual target constraints. Restoring the omitted
blocks gives a member of $\F$; appending the previously observed coordinates
gives a full record in $F$. Omitted blocks have zero uncertainty contribution.
Strong linear-programming
duality gives an optimal dual objective $\eta^*=R_B$; the batch inequalities
and the identity in \eqref{eq:random-objective} yield
$W(Q,y^*)=\sum_{g,i}w^*_{gi}v_{gi}(Q)\geq R_B$ for every $Q\in\mathcal Q_B$.
Querying the last remaining coordinate of a block, repeating a query, or
querying a coordinate already determined by the retained information adds
no constraint. Removing those redundant queries leaves an equivalent
terminal interval. Queries in omitted independent constant-coefficient
blocks can also be discarded for the width objective: their contrast
contribution is fixed, and their answers impose no constraint on other
blocks. The remaining set belongs to $\mathcal Q_B$ and costs no more.

By Proposition~\ref{prop:risk_concavity_transcript}, after fixing the seed,
the adaptive transcript has exactly the coordinate-observation fiber of
its terminal set $Q$. Its sharp width is therefore $W(Q,y)$; the query
choices and stopping decision add no constraints beyond those answers.

At $y^*$ the previous lower bound therefore applies to every terminal
transcript and seed. Taking expectations and then the worst completion
gives the lower bound $R_B$ for every adaptive policy. Randomized batches
are a subclass of adaptive policies, and the compact formulation constructs
a batch distribution with worst expected width $R_B$, proving equality.
\end{proof}

This argument uses established linear-programming duality and minimax
reasoning. Its consequence here is a checkable hardest-record certificate
for the existing structured optimizer. It does not optimize the expected
number of requests until certification, the probability that an interval
excludes zero, or the width on a particular realized historical record.

\paragraph{Extracting a completion from the implementation.}
The numerical compact implementation explicitly constrains $t_g\geq0$,
although these bounds are redundant in the mathematical primal. Its
nonnegative scenario multipliers may therefore satisfy
$\sum_i\widehat w_{gi}\leq E_g$ rather than equality. In exact arithmetic,
add each nonnegative mass deficit to a fixed coordinate. Since
$v_{gi}(Q)\geq0$, this cannot decrease any batch payoff. The resulting
completion is feasible, and its minimum batch payoff is both at least the
dual value and at most the primal value; strong duality makes these equal.
The verification code retains the raw multipliers and records clipping,
any tiny excess-mass correction, and the added deficits. It checks the
completed record with the existing weighted pricing oracle and, on small
instances, by independently enumerating every final coordinate subset
of size at most $B$, including redundant full-block queries. Numerical
discrepancies are reported separately from the exact theorem.
